\section{Robustness and Sensitivity (SS11--SS13)}
\label{sec:robustness}

The diagnostics and estimates depend on three modeling choices---the
clustering that builds the state space, the treatment of censored traces, and
the bootstrap path---and this section reports a controlled sensitivity study
for each.

\paragraph{SS11: clustering and featurization sensitivity.}
Because the certificate is a \emph{joint} test of featurization, clustering,
and the first-order assumption, its verdict could in principle be an artifact
of the clustering. We therefore re-ran the SS10 fit across $18$
configurations: two featurizations (action-category one-hot, with and without
a step-depth feature), three agglomerative linkages (Ward, average, complete),
and three cluster-count caps $k_{\max}\in\{6,8,10\}$. Table~\ref{tab:ss11}
reports the result: the composite verdict is \emph{invariant}, with all $18$
configurations rejecting the first-order absorbing DTMC. The held-out
first-passage KS test rejects ($p<10^{-3}$) in every configuration, including
the three category-only settings at $k_{\max}{=}10$ where the AIC order-test
component alone would accept a first-order chain---so the held-out KS layer is
the binding, clustering-robust component of the certificate, exactly as
intended. The recovered state count $m$ tracks $k_{\max}$ (range $6$--$10$) and
the silhouette varies with featurization, but $\hat R_\infty$ is stable across
all configurations ($0.660$--$0.674$): neither the estimate nor the verdict
moves appreciably with the clustering choice.

\emph{Reconciling the two SWE-bench $R_\infty$ estimates.} The
$\hat R_\infty$ column of Table~\ref{tab:ss11} and the $R_\infty=0.656$ of
SS14 come from different fits of the same benchmark, so each must be read
against its own target: the $160$-trace sub-sample used here has a held-out
pass rate of $0.663$, whereas SS14's full $479$-trace corpus has $0.683$. Read
that way, both are well calibrated---every configuration in the grid lands
within $0.012$ of $0.663$, and SS14 within $0.027$ of $0.683$---so neither
sub-sampling nor the choice of state space distorts the asymptotic estimate.
Because it uses the whole corpus and needs no clustering, the calibration claim
of \S\ref{sec:ss14} and the SWE-bench point of Fig.~\ref{fig:calibration} are
taken from the SS14 fit; $R_\infty$ should in general be reported against the
held-out rate of the split that produced it.

% Auto-generated by SS11_ablation.py
\begin{table}[!t]\centering
\caption{SS11: clustering and featurization ablation ($18$ configurations).}
\label{tab:ss11}\scriptsize
\setlength{\tabcolsep}{3.5pt}
\begin{tabular}{llrrrrrl}
\toprule
feat. & linkage & $k_{\max}$ & $m$ & silh. & $\Delta_{\mathrm{AIC}}$ & KS $p$ & verdict \\
\midrule
cat+depth & ward & 6 & 6 & 0.78 & +872 & 0.000 & REJECT \\
cat+depth & ward & 8 & 7 & 0.80 & +817 & 0.000 & REJECT \\
cat+depth & ward & 10 & 7 & 0.80 & +817 & 0.000 & REJECT \\
cat+depth & average & 6 & 6 & 0.78 & +872 & 0.000 & REJECT \\
cat+depth & average & 8 & 7 & 0.80 & +817 & 0.000 & REJECT \\
cat+depth & average & 10 & 7 & 0.80 & +817 & 0.000 & REJECT \\
cat+depth & complete & 6 & 6 & 0.78 & +872 & 0.000 & REJECT \\
cat+depth & complete & 8 & 7 & 0.80 & +817 & 0.000 & REJECT \\
cat+depth & complete & 10 & 7 & 0.80 & +817 & 0.000 & REJECT \\
cat & ward & 6 & 6 & 0.95 & +872 & 0.000 & REJECT \\
cat & ward & 8 & 8 & 0.99 & +736 & 0.000 & REJECT \\
cat & ward & 10 & 10 & 1.00 & -126 & 0.000 & REJECT \\
cat & average & 6 & 6 & 0.95 & +872 & 0.000 & REJECT \\
cat & average & 8 & 8 & 0.99 & +736 & 0.000 & REJECT \\
cat & average & 10 & 10 & 1.00 & -126 & 0.000 & REJECT \\
cat & complete & 6 & 6 & 0.95 & +872 & 0.000 & REJECT \\
cat & complete & 8 & 8 & 0.99 & +736 & 0.000 & REJECT \\
cat & complete & 10 & 10 & 1.00 & -126 & 0.000 & REJECT \\
\bottomrule
\end{tabular}
\end{table}


\paragraph{SS12: censoring sensitivity.}
The censoring figures reported elsewhere ($\le 10\%$) are mild, and real
deployments may truncate far more aggressively. We took a known ground-truth
absorbing chain ($R_\infty = 0.784$), sampled $600$ trajectories, and refit
\textsc{TraceToChain} after right-censoring a controlled fraction of traces
(truncating at a uniform-random transient step and marking the trace as
censored, neither success nor failure). Under this non-informative censoring the recovered
$\hat R_\infty$ stays within $\approx\!4\%$ of the truth at every level: the
relative bias moves only over $[-4.2\%, -2.6\%]$ across censoring fractions
$0$--$50\%$, against $-3.3\%$ with no censoring at all, so censoring
contributes almost none of the residual. Censored traces still contribute
their observed transient transitions but are correctly excluded from the
absorption counts, so the exit estimates are not materially biased. This
holds for censoring independent of the outcome; \emph{informative} censoring
(e.g.\ a step budget that truncates precisely the long, failure-bound
trajectories) could still bias the estimates, and detecting that regime is
left to future work.

%% [REDUCTION] tab:ss12 is a nearly-flat 7-row column whose every value is
%% quoted inline below; retained in the artifact.
%% % Auto-generated by SS12_censoring.py
\begin{table}[!t]\centering
\caption{SS12: censoring sensitivity on a ground-truth chain ($R_\infty=0.784$).}
\label{tab:ss12}
\begin{tabular}{rrrr}
\toprule
censoring & $\hat m$ & $\hat R_\infty$ & rel.\ bias \\
\midrule
0\% & 5 & 0.758 & -3.3\% \\
5\% & 5 & 0.755 & -3.7\% \\
10\% & 5 & 0.755 & -3.7\% \\
20\% & 5 & 0.764 & -2.6\% \\
30\% & 5 & 0.755 & -3.7\% \\
40\% & 5 & 0.752 & -4.2\% \\
50\% & 5 & 0.756 & -3.6\% \\
\bottomrule
\end{tabular}
\end{table}


\paragraph{SS13: bootstrap path and uncertainty understatement.}
The intervals of \S\ref{sec:uq} are reported via the fast bootstrap path, which
holds the target clustering fixed (nearest-centroid reassignment of resampled
steps). The module also implements a full re-cluster bootstrap
(\texttt{fast=False}) that re-runs Ward clustering on every resample and so
additionally captures clustering instability. On a controlled $250$-trace
corpus the full re-cluster path gives transition intervals whose median
half-width is $\approx\!4.7\times$ that of the fixed-clustering path ($0.158$
vs.\ $0.034$). The fixed-clustering intervals therefore understate total
uncertainty and should be read as a lower bound conditioned on the chosen
clustering; we recommend reporting the full re-cluster path (or both) when
clustering stability is in question.
